% Revize 2026-10-28: obsah, matematika a návaznost; podrobnosti a časy
% viz planovani/REVIZE_PRIPRAV_1_ROCNIK.md.
% Aktuální porovnání se skutečným obsahem místních PDF Realisticky.cz:
% externi_zdroje/realisticky/01_Zakladni_poznatky/02_Ciselne_obory/010201_Ciselne_obory_prirozena_cisla/010201_Lekce.pdf
% externi_zdroje/realisticky/01_Zakladni_poznatky/02_Ciselne_obory/010205_Realna_cisla_I/010205_Lekce.pdf
% externi_zdroje/realisticky/01_Zakladni_poznatky/04_Formalni_logika_vyroky/010408_Kvantifikovane_vyroky/010408_Lekce.pdf
% externi_zdroje/realisticky/01_Zakladni_poznatky/04_Formalni_logika_vyroky/010409_Negace_kvantifikovanych_vyroku/010409_Lekce.pdf
% Rozpis: 1_rocnik/tematicke_plany/Rozpis_vyuky_MAT_1_rocnik.tex, hodina 11;
% SPS.pdf, oddíl 1.2. Opakuje 02–09; připravuje výčet a vlastnost množin v 12.
% Učebnice: ucebnice/01_Zakladni_poznatky/Vyrokova_logika/Výroková logika.tex,
% Kvantifikované výroky a jejich negace, Implikace; Ciselne_obory/Číselné obory.tex,
% Druhy číselných oborů a Absolutní hodnota.
% Realisticky.cz, místní knihovna externi_zdroje/realisticky/01_Zakladni_poznatky/:
% 04_Formalni_logika_vyroky/010408_Kvantifikovane_vyroky/010408_Lekce.pdf
% (Př. 1 a 5: význam kvantifikace, pravdivost obou kvantifikací);
% 04_Formalni_logika_vyroky/010409_Negace_kvantifikovanych_vyroku/010409_Lekce.pdf
% (Př. 3–7 a tabulka);
% 04_Formalni_logika_vyroky/010410_Formalni_logika_shrnuti/010410_Lekce.pdf
% (úvodní Zádrhele).
% Inspirace postupem svědek/protipříklad, nikoli doslovné převzetí.
% Všechna konkrétní zadání a řešení jsou vlastní tvorba; chyby z testu
% nelze znát předem, níže jsou výslovně modelové situace.
\documentclass[11pt,a4paper]{article}
% Spolecny styl pro ucitelske pripravy MAT (Work/MAT)
% Verze 1.1, 2026-10-23
\NeedsTeXFormat{LaTeX2e}
\ProvidesPackage{priprava}[2026/10/23 v1.1 Sablona ucitelskych priprav MAT]

% Zakladni typografie a cestina
\RequirePackage[T1]{fontenc}
\RequirePackage[utf8]{inputenc}
\RequirePackage[czech]{babel}
\RequirePackage{lmodern}
\RequirePackage{microtype}
\RequirePackage{amsmath,amssymb}
\RequirePackage{xcolor}
\RequirePackage{enumitem}
\RequirePackage[most]{tcolorbox}
\RequirePackage{tabularx,booktabs,array}
\RequirePackage{fancyhdr}
\RequirePackage{needspace}
\RequirePackage{geometry}
\RequirePackage[hidelinks]{hyperref}

\geometry{a4paper,top=17mm,bottom=18mm,left=23mm,right=23mm}
\definecolor{notegray}{gray}{0.39}
\setlength{\parindent}{0pt}
\setlength{\parskip}{2.5pt plus 1pt minus 0.5pt}
\setlist[itemize]{topsep=2pt,itemsep=1pt,parsep=0pt,leftmargin=1.45em}
\setlist[enumerate]{topsep=2pt,itemsep=2pt,parsep=0pt,leftmargin=1.65em}
\setcounter{secnumdepth}{0}
\renewcommand{\arraystretch}{1.12}

% Zahlavi: vlevo rocnik | globalni tema, vpravo cislo: nazev hodiny
% Zapati zustava zcela prazdne.
\pagestyle{fancy}
\fancyhf{}
\renewcommand{\headrulewidth}{0.35pt}
\renewcommand{\footrulewidth}{0pt}
\setlength{\headheight}{14pt}
\newcommand{\nastavPripravu}[4]{%
  \fancyhead[L]{\small #1\enspace\textbar\enspace #2}%
  \fancyhead[R]{\small #3: #4}%
}

% Minimalisticke cernobile boxy. Kazda definice zvlast.
\tcbset{commonlesson/.style={enhanced,breakable=false,colback=white,colframe=black,
  colbacktitle=white,coltitle=black,boxrule=0.45pt,
  arc=0.8mm,outer arc=0.8mm,left=3mm,right=3mm,top=1.6mm,bottom=1.6mm,
  toptitle=1mm,bottomtitle=0.7mm,fonttitle=\bfseries,
  attach boxed title to top left={xshift=2.5mm,yshift=-1.15mm},
  boxed title style={colback=white,colframe=white,boxrule=0pt,left=1mm,right=1mm}}}
\newtcolorbox{zadaniBox}[1][]{commonlesson,title={Zadání},#1}
\newtcolorbox{prikladBox}[1][]{commonlesson,title={Příklad},#1}
\newtcolorbox{definiceBox}[1][]{commonlesson,title={Definice},#1}

% Spolecny box pro zadani a jeho reseni, bez mezery mezi castmi.
% Pouziti:
%   \begin{prikladReseniBox}
%     Zadani ulohy.
%     \tcblower
%     \begin{align*} ... \end{align*}
%   \end{prikladReseniBox}
% Spodni cast ma jemny sedy podklad a delici linku; lze v ni sazet
% vice radku matematiky. Obe casti maji jen jeden spolecny ramecek.
\tcbset{solutionlesson/.style={
  commonlesson,
  enhanced jigsaw,
  breakable,
  colbacklower=black!3,
  segmentation style={solid,draw=black!25,line width=0.35pt},
  fontlower=\small\color{notegray},
  before lower={\noindent\textbf{Řešení:}\enspace},
  leftlower=3mm,rightlower=3mm
}}
\newtcolorbox{prikladReseniBox}[1][]{solutionlesson,title={Příklad},#1}
\newtcolorbox{zadaniReseniBox}[1][]{solutionlesson,title={Zadání},#1}

% Pomucky pro autora pripravy: sedive poznamky; \reseni pro starsi pripravy
\newcommand{\poznamka}[1]{\par\begingroup\small\color{notegray}\itshape
  \textnormal{\textbf{Poznámka pro učitele.}}\enspace #1\par\endgroup}
\newcommand{\reseni}[1]{\par\begingroup\small\color{notegray}
  \textbf{Očekávané řešení:}\enspace #1\par\endgroup}
\newcommand{\blok}[1]{\needspace{5\baselineskip}\subsection*{#1}}
\newcommand{\vyrok}[1]{\uv{#1}}

\nastavPripravu{1. ročník}{Číselné obory}{11}{Číselné obory a výroky}
\begin{document}

\poznamka{Orientačně: skutečné chyby z testu 5 min, obor a kvantifikace
8 min, řešené příklady 9 min, samostatná práce s rozborem 11 min,
shrnutí 2 min; s návratem k diagramu základ asi 42–45 min.
Rezerva 10–15 min k výběru, domácí
práce 15–20 min. Modelové chyby použij jen tam, kde odpovídají potřebám
třídy; netvrď, že se v písemkách opravdu vyskytly.}

\blok{Stejná formule, jiný výběr čísel}
Vrať dvě až tři skutečné chyby z předchozího testu jako krátké opravy
postupu. Ptej se: \uv{Co je zadáno? Odkud číslo vybíráme? Chceme všechny,
nebo stačí jedno?} Nemáme-li výsledky opravené, začni následujícím příkladem.
Konvence zůstává $\mathbb N=\{1,2,3,\ldots\}$.

\blok{Návrat k diagramu: kde leží svědek nebo protipříklad?}
\begin{center}
\begin{tikzpicture}[x=1cm,y=1cm,font=\small]
\draw[rounded corners,fill=black!4] (0,0) rectangle (12,3.9);
\node at (11.4,3.5) {$\mathbb R$};
\draw[rounded corners,fill=white] (.3,.3) rectangle (8.5,3.5);
\node at (7.8,3.1) {$\mathbb Q$};
\draw[rounded corners,fill=black!8] (.6,.6) rectangle (5.6,3.1);
\node at (4.9,2.7) {$\mathbb Z$};
\draw[rounded corners,fill=white] (.9,.9) rectangle (3.2,2.7);
\node at (2.1,2.3) {$\mathbb N$};
\node at (2.1,1.5) {$1,2,3,\ldots$};
\node at (4.4,1.5) {$0,-1$};
\node at (7,1.5) {$\frac12,-\frac32$};
\node[align=center] at (10.1,2.5) {$\mathbb R\setminus\mathbb Q$};
\node at (10.1,1.5) {$\sqrt2,\pi$};
\end{tikzpicture}
\end{center}
Nech žáky umístit $\frac12$ a $\sqrt2$ a vyjmenovat všechny obory,
do nichž patří. Potom bez výpočtu přečtěte tabulku existence:
\begin{center}
\begin{tabular}{lcccc}
\toprule
Existuje $x$ v daném oboru s vlastností & $\mathbb N$ & $\mathbb Z$ & $\mathbb Q$ & $\mathbb R$\\
\midrule
$x=0$ & ne & ano & ano & ano\\
$2x=1$ & ne & ne & ano & ano\\
$x^2=2$ & ne & ne & ne & ano\\
\bottomrule
\end{tabular}
\end{center}
V posledním řádku v $\mathbb R$ stačí svědek $\sqrt2$; všechny
možnosti jsou $\sqrt2$ a $-\sqrt2$, obě iracionální.
Při rozšíření oboru se pravdivý existenční výrok se stejnou vlastností
nemůže stát nepravdivým: původní svědek zůstává k dispozici.
U obecného výroku je tomu opačně: další prvky mohou přinést
protipříklad. Například $\forall x\in\mathbb N:x>0$ platí,
ale v $\mathbb Z$ jej vyvrátí $0$.

\poznamka{Práce s diagramem a tabulkou je návratem k hodině 07,
asi na 7–10 minut. Nezavádí nové množinové operace.
Zařaď ji před P1 a vynech část opakovaného slovního shrnování.}

\begin{definiceBox}
\textbf{Obor proměnné} určuje, která čísla smíme za proměnnou dosazovat.
U výrokové formy $P(x)$ je součástí zadání; při změně oboru se může
změnit pravdivost kvantifikovaného výroku.
\end{definiceBox}

\begin{definiceBox}
Připomenutí: \textbf{obecný výrok} $\forall x\in M:\ P(x)$ říká,
že vlastnost platí pro každý prvek určeného oboru $M$.
Jediný nevyhovující prvek je protipříkladem a výrok vyvrací.
\end{definiceBox}

\begin{definiceBox}
Připomenutí: \textbf{existenční výrok} $\exists x\in M:\ P(x)$ říká,
že existuje alespoň jeden vyhovující prvek daného oboru.
Takový prvek můžeme uvést jako svědka pravdivosti.
\end{definiceBox}

\begin{prikladReseniBox}[title={P1: měníme pouze obor},breakable=false]
Rozhodni o $\exists x\in M:\ 2x=1$ postupně pro
$M=\mathbb N$, $\mathbb Z$, $\mathbb Q$ a $\mathbb R$.
Potom stejně posuď $\forall x\in M:\ x>0$.
\tcblower
Rovnost $2x=1$ vyžaduje jedinou hodnotu $x=\frac12$.
V $\mathbb N$ a $\mathbb Z$ proto existenční výrok neplatí,
v $\mathbb Q$ a $\mathbb R$ platí.

Obecný výrok $x>0$ platí v $\mathbb N$ podle definice.
V $\mathbb Z,\mathbb Q,\mathbb R$ neplatí, například kvůli $x=0$.
Polovinu nelze použít jako svědka výroku, jehož oborem je $\mathbb Z$.
\end{prikladReseniBox}

Při důkazu obecného výroku o nekonečném oboru nestačí několik pokusů.
Musíme zdůvodnit všechny možnosti. Naopak protipříklad hledáme jediný,
ale musí pocházet ze zadaného oboru.

\blok{Negace zachová obor a obrátí požadavek}
Z hodiny 04 připomeň:
\[
\neg(\forall x\in M:P(x))\ \Leftrightarrow\ \exists x\in M:\neg P(x),
\]
\[
\neg(\exists x\in M:P(x))\ \Leftrightarrow\ \forall x\in M:\neg P(x).
\]
Při negaci se obor $M$ nemění. Negace $x>0$ je $x\le0$, nikoli $x<0$.

\begin{prikladReseniBox}[title={P2: omezení oboru se projeví i v negaci},breakable=false]
Posuď $V:\ \forall x\in M:\ x\ge1$ pro $M=\mathbb N$ a pro $M=\mathbb Q$.
Pro oba obory napiš negaci a dolož její pravdivost nebo nepravdivost.
\tcblower
Pro $\mathbb N$ je $V$ pravdivý, protože každé přirozené číslo je alespoň $1$.
Negace $\exists x\in\mathbb N:\ x<1$ je nepravdivá.

Pro $\mathbb Q$ je $V$ nepravdivý, například pro $x=\frac12$.
Negace $\exists x\in\mathbb Q:\ x<1$ je pravdivá se stejným svědkem.
Negace neříká, že menší než $1$ musejí být všechna racionální čísla.
\end{prikladReseniBox}

\begin{prikladReseniBox}[title={P3: nezaměňuj implikaci s jejím obrácením},breakable=false]
Pro libovolné reálné $x$ porovnej tvrzení:
\[
x<0\Rightarrow |x|>0,\qquad |x|>0\Rightarrow x<0.
\]
\tcblower
První implikace platí pro každé reálné $x$: když $x<0$, je $|x|=-x>0$;
při nepravdivém předpokladu je implikace rovněž pravdivá.
Druhá obecně neplatí: pro $x=2$ je předpoklad pravdivý a závěr nepravdivý.
Samotná kladná vzdálenost od nuly nerozhoduje o straně osy.
\end{prikladReseniBox}

\blok{Samostatné procvičení}
\begin{zadaniReseniBox}[title={S1: pravdivost vždy dolož},breakable=false]
U každého výroku rozhodni o pravdivosti, uveď zdůvodnění nebo protipříklad
a napiš jeho negaci. Značky P a N v řešení znamenají pravdu a nepravdu.
\begin{enumerate}[label=\alph*)]
\item $\forall x\in\mathbb Z:\ x^2\ge0$.
\item $\exists x\in\mathbb Z:\ 0<x<1$.
\item $\exists x\in\mathbb Q:\ 0<x<1$.
\item $\forall x\in\mathbb R:\ |x|=x$.
\item $\forall x\in\langle-1;1\rangle:\ x<1$.
\end{enumerate}
\tcblower
a) P: nula má čtverec $0$, jinak násobíme dvě čísla stejného znaménka.
Negace $\exists x\in\mathbb Z:\ x^2<0$ je N.

b) N: mezi sousedními celými čísly $0$ a $1$ další celé není.
Negace $\forall x\in\mathbb Z:\ x\le0\ \lor\ x\ge1$ je P.

c) P: například $x=\frac12$. Negace
$\forall x\in\mathbb Q:\ x\le0\ \lor\ x\ge1$ je N.

d) N: pro $x=-1$ je $|x|=1\ne-1$.
Negace $\exists x\in\mathbb R:\ |x|\ne x$ je P.

e) N: $x=1$ patří do oboru a nesplňuje ostrou nerovnost.
Negace $\exists x\in\langle-1;1\rangle:\ x\ge1$ je P.
\end{zadaniReseniBox}

\begin{zadaniReseniBox}[title={S2: oprav modelové úvahy},breakable=false]
\begin{enumerate}[label=\alph*)]
\item \uv{Dosadil jsem $1,2,3$ a vždy vyšlo $x^2>x$. Platí to tedy pro všechna
reálná čísla.} Najdi chybu v dosazení i v zobecnění.
\item \uv{Výrok $\exists x\in\mathbb Z:\ 2x=3$ platí, protože $x=1{,}5$.}
\item \uv{Negace $x<2$ je $x>2$.} Ověř obě tvrzení pro $x=2$.
\end{enumerate}
\tcblower
a) Už pro $x=1$ je $1>1$ nepravda. Ani několik správných dosazení by
nedokázalo obecný výrok o $\mathbb R$; dalším protipříkladem je $x=0$.
b) Hodnota $1{,}5$ není celé číslo, a nelze ji tedy použít jako svědka.
Jiné řešení $2x=3$ není; existenční výrok je nepravdivý.
c) Pro $x=2$ jsou obě ostré nerovnosti nepravdivé, takže nejde o negaci.
Správně $x\ge2$.
\end{zadaniReseniBox}

\poznamka{Na konci zvol obor celých čísel a podmínku $-1\le x<3$.
Nech žáky říci všechna vyhovující čísla: $-1,0,1,2$. Po změně oboru
na reálná už nelze vypsat všechna jednotlivě. Tím připravíš hodinu 12;
formální zadávání množin a další definice ponech jí.}

\blok{Rezerva: od příkladu k obecnému zdůvodnění}
\begin{zadaniReseniBox}[title={R1: volitelně, 10–15 minut},breakable=false]
\begin{enumerate}[label=\alph*)]
\item Platí $\forall x\in\mathbb Z:\ x^2\ge x$? Dokaž nebo vyvrať.
Co se změní po nahrazení $\mathbb Z$ oborem $\mathbb Q$?
\item Pro $M=\{-2,0,3\}$ urči pravdivost
$\forall x\in M:\ (x>0\Rightarrow x\ge1)$. Napiš negaci.
Pak posuď původní výrok pro $M=\mathbb Q$.
\item Utvoř pravdivý obecný a pravdivý existenční výrok z téže vlastnosti
$x^2+1>0$ v oboru $\mathbb R$. Vysvětli rozdíl v jejich sdělení.
\end{enumerate}
\tcblower
a) V $\mathbb Z$ platí: pro $x\le0$ je $x^2\ge0\ge x$;
pro celé $x\ge1$ násobíme $x\ge1$ kladným $x$, takže $x^2\ge x$.
Jiné celé hodnoty nejsou. V $\mathbb Q$ neplatí: $x=\frac12$ dává
$\frac14<\frac12$.

b) Ve výčtu platí: pro $-2$ a $0$ je předpoklad nepravdivý,
pro $3$ je předpoklad i závěr pravdivý.
Negace $\exists x\in M:\ x>0\ \land\ x<1$ v tomto $M$ neplatí.
V $\mathbb Q$ je původní obecný výrok nepravdivý kvůli $x=\frac12$.

c) $\forall x\in\mathbb R:\ x^2+1>0$ je pravdivý, neboť $x^2+1\ge1$.
$\exists x\in\mathbb R:\ x^2+1>0$ také platí, například pro $x=0$.
První říká, že vyhovují všichni, druhý požaduje alespoň jednoho.
\end{zadaniReseniBox}

\blok{Domácí procvičování}
\begin{zadaniReseniBox}[title={D1: přibližně 15–20 minut},breakable=false]
U 1–4 urči pravdivost, zdůvodni a zapiš negaci.
\begin{enumerate}[label=\arabic*)]
\item $\forall n\in\mathbb N:\ n\ge1$.
\item $\exists x\in\mathbb Z:\ 2x=-3$.
\item $\exists x\in\mathbb Q:\ 2x=-3$.
\item $\forall x\in(-1;1):\ |x|<1$.
\item Posuď pro všechny reálné hodnoty tvrzení
$x>2\Rightarrow x>0$ a jeho obrácení. Nepravdivé tvrzení vyvrať.
\item Vyber všechna celá čísla s $-2<x\le2$. Rozhodni o pravdivosti
výroku \uv{Každé z vybraných čísel je kladné}. Utvoř negaci.
\end{enumerate}
\tcblower
1. P podle definice $\mathbb N$; negace $\exists n\in\mathbb N:\ n<1$ je N.
2. N, jediné možné $x=-\frac32$ není celé;
negace $\forall x\in\mathbb Z:\ 2x\ne-3$ je P.
3. P se svědkem $x=-\frac32$; negace $\forall x\in\mathbb Q:\ 2x\ne-3$ je N.
4. P: každý bod otevřeného úseku mezi $-1$ a $1$ je od nuly vzdálen
méně než $1$; negace $\exists x\in(-1;1):\ |x|\ge1$ je N.
5. První platí, protože $x>2>0$; obrácení neplatí například pro $x=1$.
6. $-1,0,1,2$; obecný výrok je N. Negace: alespoň jedno z vybraných
čísel je nekladné; je P, například díky $-1$ nebo $0$.
\end{zadaniReseniBox}

\end{document}
